% Capacidades de texforge: documento de referencia editorial.
% Compilado con texforge (motor tectonic embebido). Leer en UTF-8.
\documentclass[11pt,a4paper]{article}

\usepackage[spanish]{babel}
\usepackage[T1]{fontenc}
\usepackage[lining]{ebgaramond}
\usepackage[scale=0.92]{sourcesanspro}
\usepackage{microtype}
\usepackage[a4paper,left=27mm,right=24mm,top=30mm,bottom=32mm,
  headheight=14pt,headsep=14mm,footskip=14mm]{geometry}
\usepackage{xcolor}
\usepackage{enumitem}
\usepackage{fancyhdr}
\usepackage{titlesec}
\usepackage{tocloft}
\usepackage{needspace}
\usepackage{etoolbox}
\usepackage{booktabs}
\usepackage{array}
\usepackage{tabularx}
\usepackage{float}
\usepackage{placeins}
% Las figuras del capitulo de diagramas son altas: dejan que una figura
% ocupe casi toda la pagina y que quepan unas lineas de texto a su lado,
% en vez de empujarla a una pagina propia.
\renewcommand{\topfraction}{0.92}
\renewcommand{\bottomfraction}{0.85}
\renewcommand{\textfraction}{0.06}
\renewcommand{\floatpagefraction}{0.85}
% Las figuras del capitulo de diagramas van seguidas unas de otras: el aire
% por defecto entre figura y texto deja huecos que empujan una figura a la
% pagina siguiente.
\setlength{\textfloatsep}{11pt plus 2pt minus 2pt}
\setlength{\intextsep}{9pt plus 2pt minus 2pt}
\setlength{\floatsep}{11pt plus 2pt minus 2pt}
\setlength{\abovecaptionskip}{6pt}
\usepackage{graphicx}
\usepackage{lastpage}
\usepackage{tikz}
\usepackage{pgfplots}
\pgfplotsset{compat=1.18}
\usepackage[most]{tcolorbox}
\usepackage{amsmath,amssymb}
\usepackage[hyphens]{url}
\usepackage{hyperref}

% Paleta definida una sola vez: tonos del laboratorio y un unico acento.
\definecolor{benchparchment}{HTML}{F4F0EA}
\definecolor{benchdeep}{HTML}{EEE7DC}
\definecolor{benchlamp}{HTML}{F6ECD2}
\definecolor{benchink}{HTML}{1D1510}
\definecolor{benchsepia}{HTML}{6B5A43}
\definecolor{benchgraphite}{HTML}{5D5852}
\definecolor{benchbrass}{HTML}{8A6A2C}
\definecolor{benchcopper}{HTML}{9A5530}
\definecolor{benchgarnet}{HTML}{87222C}
\definecolor{benchslate}{HTML}{222928}
\colorlet{docink}{benchink}
\definecolor{docaccent}{HTML}{87222C}
\colorlet{docmuted}{benchsepia}
\colorlet{docrule}{benchsepia!35!white}
\colorlet{docwarn}{benchgarnet}
\hypersetup{colorlinks=true,linkcolor=docink,urlcolor=docaccent,
  citecolor=docaccent,breaklinks=true,
  pdftitle={Capacidades de texforge},pdfauthor={Jheison M}}
\urlstyle{sf}

% Etiquetas fijas en espanol.
\newcommand{\lblVersion}{Versión}
\newcommand{\lblDate}{Fecha}
\newcommand{\lblAuthor}{Autor}
\newcommand{\lblOrganization}{Organización}
\newcommand{\lblNote}{Nota}
\newcommand{\lblNews}{Novedad 1.0}
\newcommand{\lblPage}{Página}
\newcommand{\lblOf}{de}

% Metadatos del documento, usados por la portada y el encabezado.
\newcommand{\DocTitle}{Capacidades de texforge}
\newcommand{\DocSubtitle}{Lo que el compilador hace por tus documentos, contado en un documento hecho con él}
\newcommand{\DocOrg}{UniverLab}
\newcommand{\DocOwner}{Jheison M}
\newcommand{\DocVersion}{Versión 1.0}
\newcommand{\DocDate}{18 de noviembre de 2026}

% Sin soporte para xetex: las versalitas llevan el enfasis.
\renewcommand{\textls}[2][]{#2}
\newcommand{\caps}[1]{{\sffamily\scshape #1}}
\newcommand{\metalabel}[1]{{\sffamily\scriptsize\scshape\textls[90]{\color{docmuted}#1}}}

% Cabecera y pie: organizacion y titulo arriba, pagina abajo.
\fancypagestyle{doc}{%
  \fancyhf{}
  \fancyhead[L]{\small\caps{\DocOrg}}
  \fancyhead[R]{\small\itshape\DocTitle}
  \fancyfoot[L]{\small\caps{\DocVersion}}
  \fancyfoot[R]{\small\sffamily\lblPage\ \thepage\ \lblOf\ \pageref{LastPage}}
  \renewcommand{\headrulewidth}{0.4pt}
  \renewcommand{\footrulewidth}{0.4pt}
  \renewcommand{\headrule}{\hbox to\headwidth{\color{docrule}\leaders\hrule height \headrulewidth\hfill}}
  \renewcommand{\footrule}{\hbox to\headwidth{\color{docrule}\leaders\hrule height \footrulewidth\hfill}}}
\pagestyle{doc}
\fancypagestyle{plain}{\pagestyle{doc}}

% Titulos: seccion Garamond con numero en acento y regla, subseccion sans.
\titleformat{\section}{\normalfont\fontsize{20}{24}\selectfont\color{docink}}
{\textcolor{docaccent}{\thesection}}{0.7em}{}[{\vspace{3pt}\color{docrule}\titlerule}]
\titlespacing*{\section}{0pt}{1.6\baselineskip}{0.7\baselineskip}
\titleformat{\subsection}{\sffamily\bfseries\normalsize\color{docink}}
{\textcolor{docaccent}{\thesubsection}}{0.7em}{}
\titlespacing*{\subsection}{0pt}{1.0\baselineskip}{0.3\baselineskip}
\titleformat{\subsubsection}{\normalfont\itshape\color{docink}}{}{0pt}{}
\titlespacing*{\subsubsection}{0pt}{0.8\baselineskip}{0.2\baselineskip}

% Indice con tocloft.
\renewcommand{\cftsecfont}{\normalfont\bfseries}
\renewcommand{\cftsecpagefont}{\normalfont\bfseries}
\renewcommand{\cftsubsecfont}{\normalfont}
\renewcommand{\cftsecdotsep}{\cftnodots}
\renewcommand{\cftsubsecdotsep}{\cftdotsep}
\setlength{\cftbeforesecskip}{0.5em}
\setlength{\cftbeforesubsecskip}{0.1em}
\renewcommand{\cftdotsep}{2}
\renewcommand{\cfttoctitlefont}{\normalfont\fontsize{20}{24}\selectfont\color{docink}}
\setlength{\cftaftertoctitleskip}{0.8em}

\frenchspacing
\let\origsection\section
\renewcommand{\section}{\needspace{8\baselineskip}\origsection}
\setlength{\parindent}{0pt}
\setlength{\parskip}{0.4\baselineskip}
\setlength{\emergencystretch}{2em}
\renewcommand{\arraystretch}{1.25}
\setlist{leftmargin=1.5em,itemsep=1pt,topsep=2pt,parsep=0pt}
\setlist[itemize]{label={\color{docaccent}\textbullet}}

% Longitudes de la portada: definidas aqui para que el cuerpo solo nombre comandos.
\newlength{\covertopskip}\setlength{\covertopskip}{2mm}
\newlength{\covermidskip}\setlength{\covermidskip}{18pt}
\newlength{\coversmallskip}\setlength{\coversmallskip}{2pt}

% Franja de la portada, como las plantillas: texto en el preambulo.
\newcommand{\DocKicker}{Documento de referencia \textperiodcentered\ Versión 1.0}
\newcommand{\coverbands}{%
  \begin{tikzpicture}[remember picture,overlay]
    \node[anchor=north west,inner sep=0pt,fill=benchgarnet,minimum width=\paperwidth,
      minimum height=2.4mm] at ([yshift=-22mm]current page.north west) {};
    \node[anchor=north west,inner sep=0pt,fill=benchgarnet,minimum width=\paperwidth,
      minimum height=0.9mm] at ([yshift=-132mm]current page.north west) {};
    \node[anchor=north west,inner sep=0pt,fill=benchgarnet,minimum width=\paperwidth,
      minimum height=0.35mm] at ([yshift=-134mm]current page.north west) {};
  \end{tikzpicture}}

% Cajas: panel tintado con barra de acento y rotulo en versalitas.
\tcbset{docbox/.style={enhanced,breakable,colback=#1!6,colframe=#1!6,
    boxrule=0pt,arc=0pt,outer arc=0pt,borderline west={2.4pt}{0pt}{#1},
    left=11pt,right=10pt,top=6pt,bottom=6pt,before skip=7pt,after skip=7pt,
    fonttitle=\sffamily\scriptsize\scshape\bfseries,coltitle=#1,
    attach boxed title to top left={xshift=-1pt,yshift=-1pt},
    boxed title style={empty,left=0pt,top=0pt,bottom=0pt}}}
\newtcolorbox{notebox}[1][]{docbox=docmuted,title={\textls[90]{\lblNote}},#1}
\newtcolorbox{novabox}[1][]{docbox=docaccent,title={\textls[90]{\lblNews}},#1}

\begin{document}

  % Portada: titulo y datos en un bloque enmarcado por reglas granate.
  \begingroup
  \thispagestyle{empty}
  \noindent\begin{minipage}[t][\dimexpr\textheight-2\baselineskip\relax][t]{\textwidth}
    \coverbands\par\nointerlineskip
    \raggedright\color{docink}
    \vspace*{\covertopskip}
    {\large\color{docmuted}\caps{\DocOrg}\par}
    \vspace{18mm}
    {\sffamily\small\scshape\color{docaccent}\DocKicker\par}
    \vspace{4mm}
    {\fontsize{38}{42}\selectfont\DocTitle\par}
    \vspace{6mm}
    {\fontsize{16}{20}\selectfont\itshape\color{docmuted}\DocSubtitle\par}
    \vspace{42mm}\color{docink}
    \begin{tabular}{@{}p{5.6cm}p{5.6cm}@{}}
      \metalabel{\lblVersion} & \metalabel{\lblDate} \\[\coversmallskip]
      \DocVersion & \DocDate \\[\covermidskip]
      \metalabel{\lblAuthor} & \metalabel{\lblOrganization} \\[\coversmallskip]
      \DocOwner & \DocOrg \\
    \end{tabular}
    \vfill
    {\color{docmuted}\small Compilado con texforge, sin TeX Live ni MiKTeX\par}
  \end{minipage}
  \endgroup
  \clearpage

  \section*{De un vistazo}
  \begin{tcolorbox}[docbox=docaccent,title={\textls[90]{De un vistazo}}]
    Este documento demuestra todo lo que texforge sabe componer: diagramas
    Mermaid, Graphviz y D2 generados durante la compilación, un esquema
    TikZ dibujado a mano, matemáticas con referencias cruzadas, listados de
    código en cuatro lenguajes con el estilo \texttt{ink} y numeración que
    no se copia, citas BibTeX, dos tablas \texttt{booktabs} y listas con
    propósito. Todo lo que ves salió de \texttt{texforge build}.
  \end{tcolorbox}

  \tableofcontents
  \clearpage

  % Capitulo 1: que es texforge y como se produjo este documento.
\section{Qué es texforge}\label{sec:que-es}

texforge es un compilador de documentos \LaTeX{} que se distribuye como un
único binario de línea de comandos y lleva dentro su propio motor de
composición, tectonic~\cite{tectonic}. Quien escribe no instala TeX Live ni
MiKTeX, no configura espejos de paquetes ni espera descargas de varios
gibas: crea un proyecto, escribe y compila. El motor gestionado se descarga
solo la primera vez que se compila y vive bajo \texttt{\$HOME/.texforge},
fuera del documento y fuera del control de versiones del usuario.

El flujo cabe en tres verbos. \texttt{texforge new} crea un proyecto con su
\texttt{project.toml}; \texttt{texforge check} revisa el fuente sin compilar
y avisa de citas rotas, etiquetas sin resolver y palabras desconocidas;
\texttt{texforge build} compone el PDF en \texttt{build/}. Entre la
escritura y el PDF hay además un preprocesador que detecta los entornos de
diagramas y de código, los dibuja o resalta en Rust puro y trabaja sobre
copias en el directorio de compilación: los \texttt{.tex} originales nunca
se modifican, así que lo que se versiona es siempre lo que se escribió.

El motor gestionado merece una línea propia. Tectonic vive bajo
\texttt{\$HOME/.texforge}, fuera de los proyectos, y \texttt{texforge
  doctor} lo mantiene al día: actualiza el motor gestionado a la versión
soportada y avisa si el Tectonic del \texttt{PATH} es más viejo que el que
texforge espera. Nada de esto toca la configuración del usuario ni sus
cachés personales: la compilación usa sus propios temporales y el proyecto
solo guarda el fuente, el \texttt{project.toml} y el PDF.

\begin{notebox}
  Si vienes de una distribución clásica, la diferencia se nota el primer
  día: no hay que elegir qué colección instalar ni qué espejo usar.
  El primer \texttt{build} tarda lo que tarde en llegar el motor; los
  siguientes solo componen.
\end{notebox}

Este mismo documento es la prueba. Su fuente vive en
\texttt{examples/texforge-capabilities} y su organización es la de todo
proyecto: un \texttt{main.tex} con el preámbulo editorial, nueve capítulos
en \texttt{sections/} incluidos desde el principal, la bibliografía en
\texttt{bib/references.bib} y un guion de análisis en
\texttt{scripts/stats.py} cuyos fragmentos se citan desde los listados.
Cada directorio tiene un oficio y el principal los cose: el texto en las
secciones, los datos en la bibliografía y el código de ejemplo junto al
documento que lo explica. Se compiló con \texttt{texforge build} desde
ese directorio, se revisó con \texttt{texforge check} y
\texttt{texforge spell}, se formateó con \texttt{texforge fmt} y se
inspeccionó con \texttt{texforge pdf text}, \texttt{info}, \texttt{pages} y
\texttt{check}. Nada de lo que sigue está pegado como imagen: cada
diagrama se generó al compilar, cada listado se resaltó al compilar y
cada cita se resolvió desde el \texttt{.bib}.

La portada, la cabecera con el título, el pie con la página, los
titulares con número en color de acento y las cajas de aviso salen del
preámbulo, que sigue el lenguaje visual de las plantillas de UniverLab:
Garamond para el texto, Source Sans Pro para los titulares, microtipografía
y una paleta definida una sola vez con un único acento granate. El
capítulo~\ref{sec:flujo} detalla el flujo y los comandos, el
capítulo~\ref{sec:diagramas} muestra los cuatro diagramas con sus
referencias, y el capítulo~\ref{sec:novedades} resume lo nuevo de la
versión 1.0.

  % Capitulo 2: flujo de trabajo y comandos, con la tabla de comandos.
\section{El flujo de trabajo}\label{sec:flujo}

Un proyecto texforge es un directorio con un \texttt{project.toml} que nombra
la entrada y la bibliografía. Se empieza con \texttt{texforge new
  mi-documento}, que genera el esqueleto, o con \texttt{texforge init} dentro
de un proyecto \LaTeX{} existente, que detecta el \texttt{main.tex} y la
bibliografía y escribe el \texttt{project.toml} por ti. A partir de ahí
el ciclo es corto: escribir, revisar con \texttt{check}, formatear con
\texttt{fmt} y componer con \texttt{build}. \texttt{texforge clean} retira el
directorio \texttt{build/} cuando quieres empezar la composición desde
cero, y \texttt{texforge outline} muestra el árbol de secciones del
documento para comprobar la estructura sin abrir el PDF.

\subsection{Inspeccionar el resultado}

Componer no es el último paso: el PDF también se lee desde la terminal.
\texttt{texforge pdf text} vuelca el texto tal como lo vería un lector de
pantalla o un sistema de selección de personal, con las ligaduras
normalizadas y las palabras partidas al final de línea ya reunidas, así
que es la prueba real de accesibilidad del documento. \texttt{texforge pdf
  info} informa de las páginas, de las fuentes y de si van embebidas, y de
los metadatos que viajarán con el archivo. \texttt{texforge pdf pages}
dice qué sección abre cada página en un formato pensado para comparar
compilaciones línea a línea, y \texttt{texforge pdf check} comprueba que
cada palabra significativa del fuente haya sobrevivido a la composición:
si una ligadura o una codificación se come una palabra, el aviso aparece
aquí y no en la imprenta. Para mirar el resultado con los ojos,
\texttt{texforge preview} convierte el PDF en imágenes sin salir de la
terminal. Y cuando el corrector duda de una palabra, \texttt{texforge spell}
decide: las palabras personales, como apellidos o tecnologías, se aceptan
una vez para todos los proyectos, y las de un documento concreto se guardan
con \texttt{-{}-local} en el propio proyecto, en \texttt{.texforge/spell-words}.
La lista de este documento, sin ir más lejos, guarda los nombres de las
herramientas y los términos técnicos que el diccionario no trae, para que
\texttt{check} revise prosa y no vocabulario.

\subsection{Los comandos}

La tabla~\ref{tab:comandos} resume los comandos, incluidos los que llegaron
con la versión 1.0. Las novedades de la versión se describen con calma en
el capítulo~\ref{sec:novedades}; aquí basta saber que el flujo diario
cabe en \texttt{new}, \texttt{build}, \texttt{check} y \texttt{fmt}, y que el
resto aparece cuando el documento crece: plantillas, ortografía,
diagnóstico y lectura del PDF.

\begin{table}[H]
  \centering
  \caption{Comandos de texforge, con las novedades de la versión 1.0}
  \label{tab:comandos}
  \begin{tabularx}{\linewidth}{@{}lX@{}}
    \toprule
    \sffamily\small\bfseries Comando & \sffamily\small\bfseries Qué hace \\
    \midrule
    \texttt{new -t -{}-values} & Crea un proyecto rellenando la plantilla desde un fichero de valores, que queda guardado como \texttt{values.toml}. \\
    \texttt{template add} & Añade una carpeta local como plantilla, con \texttt{-{}-force} para reemplazarla. \\
    \texttt{skill} & Muestra la destreza de agente que viaja dentro del binario. \\
    \texttt{doctor} & Diagnostica el entorno y actualiza el Tectonic gestionado; avisa si el de \texttt{PATH} es más viejo. \\
    \texttt{update} & Actualiza texforge verificando la firma minisign del paquete. \\
    \texttt{outline} & Muestra el árbol de secciones del documento. \\
    \texttt{pdf text} & Imprime el texto del PDF como lo vería un lector o un sistema de selección. \\
    \texttt{pdf info} & Informa de páginas, fuentes embebidas y metadatos. \\
    \texttt{pdf pages} & Dice qué sección abre cada página, en un formato apto para comparar compilaciones. \\
    \texttt{pdf check} & Comprueba que las palabras del fuente sobrevivieron a la composición. \\
    \texttt{build} & Compone \texttt{main.tex} a PDF en \texttt{build/}. \\
    \texttt{check} & Revisa citas, etiquetas y ortografía sin compilar. \\
    \texttt{fmt} & Formatea los \texttt{.tex} y los \texttt{.bib}. \\
    \texttt{init} & Migra un proyecto \LaTeX{} existente o guía uno nuevo. \\
    \texttt{clean} & Retira el directorio \texttt{build/}. \\
    \texttt{stats} & Cuenta palabras por sección o por archivo. \\
    \texttt{preview} & Convierte el PDF en imágenes para revisarlo. \\
    \texttt{config} & Lee y escribe la configuración global. \\
    \bottomrule
  \end{tabularx}
\end{table}

\begin{novabox}
  \texttt{outline} y los cuatro subcomandos de \texttt{pdf} son nuevos en la
  1.0: antes la estructura y el PDF resultante solo se veían abriendo el
  PDF; ahora se inspeccionan desde la terminal.
\end{novabox}

Revisar antes de componer ahorra tiempo. \texttt{texforge check} encuentra
archivos \texttt{\textbackslash input} que faltan, imágenes ausentes,
citas sin entrada en el \texttt{.bib} y pares \texttt{\textbackslash ref} y
\texttt{\textbackslash label} rotos, incluidos los \texttt{label=} de los
diagramas. El corrector ortográfico respeta el idioma declarado con
\texttt{babel} y acepta el vocabulario propio con \texttt{texforge spell}:
las palabras personales viven en el diccionario global y las de un documento
concreto se guardan con \texttt{-{}-local} en el proyecto. \texttt{texforge
  fmt} deja la indentación y los \texttt{.bib} en forma canónica, y
\texttt{fmt -{}-check} lo verifica en integración continua sin tocar nada.

  % Capitulo 3: los cuatro diagramas, cada uno con pie, etiqueta y referencia.
\section{Diagramas}\label{sec:diagramas}

texforge compone diagramas descritos como texto dentro del propio fuente,
sin Inkscape, sin Node y sin binarios externos: al compilar, cada entorno se
compone en Rust puro y el resultado entra al PDF como una figura más.
Los tres motores entienden las mismas opciones --- \texttt{width},
\texttt{pos}, \texttt{caption} y \texttt{label} --- y el estilo visual se
fija una vez en \texttt{project.toml} bajo \texttt{[diagrams]}, aquí el
preajuste \texttt{editorial}, de modo que los tres dibujos comparten tinta,
acento y aire aunque los dibujen motores distintos; la familia tipográfica
llega hasta donde cada motor la deja elegir.

\subsection{Mermaid}

La figura~\ref{fig:secuencia} muestra la secuencia de una compilación: el
usuario pide el PDF, texforge valida el fuente y, si hay diagramas, los
envía al motor de composición dentro de un bucle antes de producir el
resultado. El lenguaje es Mermaid~\cite{mermaid}, el de los diagramas de
secuencia y de flujo que ya se escriben en la documentación técnica: los
participantes se declaran con una palabra, los mensajes con flechas y el
texto de cada flecha viaja sobre ella sin sintaxis adicional.

\begin{mermaid}[width=0.9\linewidth, caption={Secuencia de compilación con texforge}, label={fig:secuencia}, pos=H]
  sequenceDiagram
  autonumber
  actor Usuario
  participant texforge
  participant Validador
  participant Motor
  Usuario->>texforge: texforge build
  activate texforge
  texforge->>Validador: analiza el fuente
  activate Validador
  Validador-->>texforge: fuente válido
  deactivate Validador
  alt hay diagramas
  texforge->>Motor: prepara el motor
  activate Motor
  loop cada diagrama
  texforge->>Motor: renderiza diagrama
  Motor-->>texforge: figura lista
  end
  Motor-->>texforge: motor libre
  deactivate Motor
  else sin diagramas
  Note over texforge,Motor: compila directo
  end
  texforge->>Motor: compone el PDF
  activate Motor
  Motor-->>texforge: PDF final
  deactivate Motor
  texforge-->>Usuario: build/ listo
  deactivate texforge
\end{mermaid}

\subsubsection{El mismo diagrama en los cuatro preajustes}

El atributo \texttt{style} aplica uno de los cuatro preajustes; la misma
secuencia se repite abajo, figuras \ref{fig:secuencia-editorial},
\ref{fig:secuencia-monochrome}, \ref{fig:secuencia-technical} y
\ref{fig:secuencia-default}.

\begin{mermaid}[style=editorial, width=0.62\linewidth, caption={\texttt{editorial}: paleta sobria con un solo acento}, label={fig:secuencia-editorial}, pos=t]
  sequenceDiagram
  autonumber
  actor Usuario
  participant texforge
  participant Validador
  participant Motor
  Usuario->>texforge: texforge build
  activate texforge
  texforge->>Validador: analiza el fuente
  activate Validador
  Validador-->>texforge: fuente válido
  deactivate Validador
  alt hay diagramas
  texforge->>Motor: prepara el motor
  activate Motor
  loop cada diagrama
  texforge->>Motor: renderiza diagrama
  Motor-->>texforge: figura lista
  end
  Motor-->>texforge: motor libre
  deactivate Motor
  else sin diagramas
  Note over texforge,Motor: compila directo
  end
  texforge->>Motor: compone el PDF
  activate Motor
  Motor-->>texforge: PDF final
  deactivate Motor
  texforge-->>Usuario: build/ listo
  deactivate texforge
\end{mermaid}

\begin{mermaid}[style=monochrome, width=0.62\linewidth, caption={\texttt{monochrome}: solo grises, para blanco y negro}, label={fig:secuencia-monochrome}, pos=t]
  sequenceDiagram
  autonumber
  actor Usuario
  participant texforge
  participant Validador
  participant Motor
  Usuario->>texforge: texforge build
  activate texforge
  texforge->>Validador: analiza el fuente
  activate Validador
  Validador-->>texforge: fuente válido
  deactivate Validador
  alt hay diagramas
  texforge->>Motor: prepara el motor
  activate Motor
  loop cada diagrama
  texforge->>Motor: renderiza diagrama
  Motor-->>texforge: figura lista
  end
  Motor-->>texforge: motor libre
  deactivate Motor
  else sin diagramas
  Note over texforge,Motor: compila directo
  end
  texforge->>Motor: compone el PDF
  activate Motor
  Motor-->>texforge: PDF final
  deactivate Motor
  texforge-->>Usuario: build/ listo
  deactivate texforge
\end{mermaid}

\begin{mermaid}[style=technical, width=0.62\linewidth, caption={\texttt{technical}: trazo uniforme, sin rellenos}, label={fig:secuencia-technical}, pos=t]
  sequenceDiagram
  autonumber
  actor Usuario
  participant texforge
  participant Validador
  participant Motor
  Usuario->>texforge: texforge build
  activate texforge
  texforge->>Validador: analiza el fuente
  activate Validador
  Validador-->>texforge: fuente válido
  deactivate Validador
  alt hay diagramas
  texforge->>Motor: prepara el motor
  activate Motor
  loop cada diagrama
  texforge->>Motor: renderiza diagrama
  Motor-->>texforge: figura lista
  end
  Motor-->>texforge: motor libre
  deactivate Motor
  else sin diagramas
  Note over texforge,Motor: compila directo
  end
  texforge->>Motor: compone el PDF
  activate Motor
  Motor-->>texforge: PDF final
  deactivate Motor
  texforge-->>Usuario: build/ listo
  deactivate texforge
\end{mermaid}

\begin{mermaid}[style=default, width=0.62\linewidth, caption={\texttt{default}: el aspecto nativo del motor}, label={fig:secuencia-default}, pos=t]
  sequenceDiagram
  autonumber
  actor Usuario
  participant texforge
  participant Validador
  participant Motor
  Usuario->>texforge: texforge build
  activate texforge
  texforge->>Validador: analiza el fuente
  activate Validador
  Validador-->>texforge: fuente válido
  deactivate Validador
  alt hay diagramas
  texforge->>Motor: prepara el motor
  activate Motor
  loop cada diagrama
  texforge->>Motor: renderiza diagrama
  Motor-->>texforge: figura lista
  end
  Motor-->>texforge: motor libre
  deactivate Motor
  else sin diagramas
  Note over texforge,Motor: compila directo
  end
  texforge->>Motor: compone el PDF
  activate Motor
  Motor-->>texforge: PDF final
  deactivate Motor
  texforge-->>Usuario: build/ listo
  deactivate texforge
\end{mermaid}

El preajuste \texttt{editorial} sirve para documentos que se leen en pantalla
o se imprimen a color y buscan un acabado sobrio; \texttt{monochrome} dibuja
solo en grises y es el indicado cuando el documento se imprime en blanco y
negro; \texttt{technical} viste los diagramas como planos o esquemas, sin
rellenos y con etiquetas de ancho fijo. El valor \texttt{default} conserva el
aspecto nativo de cada motor para no cambiar ningún documento existente. El
preajuste se fija una vez en \texttt{project.toml}, bajo \texttt{[diagrams]}
con \texttt{style = "..."}, y cada entorno puede cambiarlo con su propio
atributo \texttt{style}; omitir ambos equivale a \texttt{default}. Este
documento usa \texttt{editorial} en todo su recorrido.
\FloatBarrier

\subsection{Graphviz}

La figura~\ref{fig:arquitectura} dibuja la arquitectura como un grafo
dirigido con Graphviz~\cite{graphviz}, organizada en tres etapas: la interfaz
de línea de comandos, el preprocesador con sus tres motores y la salida que
produce el PDF. El motor de composición es Rust puro: no dibuja cajas de
\texttt{subgraph cluster}, así que las etapas se distinguen por el relleno de
sus nodos y la leyenda, no por un recuadro. Cada nodo nombra una pieza y las
flechas marcan dependencias.

\begin{graphviz}[width=1.0\linewidth, caption={Arquitectura de texforge: interfaz, preprocesador y salida}, label={fig:arquitectura}, pos=H]
  digraph Arquitectura {
    rankdir=LR
    node [style=filled, fillcolor=linen, color=sienna, fontsize=14]
    edge [color=sienna, fontsize=12]

    subgraph front {
      CLI [label="texforge CLI", shape=box, fillcolor=antiquewhite]
    }

    subgraph pre {
      Mermaid [label="Mermaid", fillcolor=cornsilk]
      Graphviz [label="Graphviz", fillcolor=cornsilk]
      D2 [label="D2", fillcolor=cornsilk]
      { rank=same; Mermaid; Graphviz; D2; }
    }

    subgraph back {
      Tectonic [label="tectonic", shape=box, fillcolor=darkseagreen]
      PDF [label="PDF", shape=doublecircle, fillcolor=white]
      { rank=same; Tectonic; PDF; }
    }

    CLI -> Mermaid [label="fuente", penwidth=2]
    CLI -> Graphviz [label="fuente", penwidth=2]
    CLI -> D2 [label="fuente", penwidth=2]
    Mermaid -> Tectonic [style=dashed, color=firebrick]
    Graphviz -> Tectonic [style=dashed, color=firebrick]
    D2 -> Tectonic [style=dashed, color=firebrick]
    Tectonic -> PDF [penwidth=2, color=firebrick, label="genera"]

    subgraph legend {
      E1 [label="origen: CLI", shape=box, fillcolor=antiquewhite]
      E2 [label="motor", shape=box, fillcolor=cornsilk]
      E3 [label="salida", shape=box, fillcolor=darkseagreen]
      E1 -> E2 [label="flujo", penwidth=2]
      E2 -> E3 [label="genera", style=dashed, color=firebrick]
    }
  }
\end{graphviz}
\FloatBarrier

\subsection{D2}

La figura~\ref{fig:pipeline} resume el trabajo como una tubería con la
sintaxis de D2: escribir, revisar, componer e inspeccionar, cada etapa con
su comando dentro de la caja, y el artefacto que produce la salida. La
arista a trazos marca el paso de verificación.

\begin{d2}[width=1.0\linewidth, caption={Tubería de trabajo, de la escritura a la inspección}, label={fig:pipeline}, pos=H]
  direction: down

  classes: {
    etapa: {
      style.border-radius: 4
    }
    control: {
      style.stroke-dash: 3
    }
  }

  autor: Autor {
    shape: person
    label: "Autor"
    class: etapa
  }

  flujo: Tubería {
    grid-columns: 2
    escribir: "Escribir\nfuente .tex" {
      shape: document
      class: etapa
    }
    revisar: "Revisar\ntexforge check" {
      shape: rectangle
      class: etapa
    }
    componer: "Componer\ntexforge build" {
      shape: hexagon
      class: etapa
    }
    inspeccionar: "Inspeccionar\ntexforge pdf info" {
      shape: rectangle
      class: etapa
    }
  }

  salida: Resultado {
    pdf: "build/\nPDF final" {
      shape: page
      class: etapa
    }
  }

  autor -> flujo.escribir: redacta
  flujo.escribir -> flujo.revisar {
    class: control
  }
  flujo.revisar -> flujo.componer
  flujo.componer -> flujo.inspeccionar
  flujo.inspeccionar -> salida.pdf: entrega

  leyenda: Leyenda {
    a: paso
    b: verificación
    a -> b: acción
    b -> a: "control" {
      class: control
    }
  }
\end{d2}
\FloatBarrier

\begin{notebox}
  Los tres entornos aceptan \texttt{label=} con o sin \texttt{caption}, y
  \texttt{texforge check} resuelve \texttt{\textbackslash ref} contra esas
  etiquetas. Si una figura pierde su etiqueta, el error aparece antes de
  compilar, no como una referencia sin resolver en el PDF.
\end{notebox}

\subsection{TikZ}

Para lo que ningún motor describe, el TikZ de siempre sigue disponible. La
figura~\ref{fig:pagina} dibuja la curva $f(x)=e^{-x^{2}/2}$ con su tangente
en $x=1$: el área bajo la curva queda sombreada, el punto de tangencia y los
ejes van anotados, y la paleta del documento (\texttt{docaccent} y los tonos
\texttt{bench}) mantiene el mismo acabado que el resto del capítulo. El
código fuente se muestra en el listado~\ref{lst:tikz-curva} para que pueda
copiarse tal cual.

\begin{figure}[H]
  \centering
  \begin{tikzpicture}
    \begin{axis}[
        width=0.82\linewidth,
        height=4.9cm,
        axis lines=middle,
        axis line style={benchsepia},
        xlabel={$x$},
        ylabel={$f(x)$},
        xmin=-2.6, xmax=2.6,
        ymin=-0.15, ymax=1.25,
        xtick={-2,-1,0,1,2},
        ytick={0.5,1},
        tick label style={font=\small\sffamily, color=docink},
        label style={font=\small\sffamily, color=docink},
        samples=120,
        domain=-2.5:2.5,
      ]
      \addplot[draw=none, fill=benchlamp, domain=-2.5:2.5] {exp(-x^2/2)} \closedcycle;
      \addplot[thick, color=docaccent] {exp(-x^2/2)};
      \addplot[dashed, color=benchsepia, domain=-0.2:2.2] {exp(-0.5) - exp(-0.5)*(x-1)};
      \addplot[only marks, mark=*, mark size=2pt, color=benchink] coordinates {(1,{exp(-0.5)})};
      \node[anchor=south west, font=\small\sffamily, color=docink] at (axis cs:1.05,0.62) {tangente en $x=1$};
      \node[anchor=south, font=\small\sffamily, color=benchsepia] at (axis cs:1.35,0.95) {$f(x)=e^{-x^{2}/2}$};
    \end{axis}
  \end{tikzpicture}
  \caption{Curva gaussiana con sombreado y tangente dibujada con TikZ}
  \label{fig:pagina}
\end{figure}

\begin{code}[lang=latex, numbers=true, caption={Código fuente de la figura~\ref{fig:pagina}}, label={lst:tikz-curva}]
\begin{tikzpicture}
  \begin{axis}[
    width=0.82\linewidth, height=4.9cm,
    axis lines=middle, axis line style={benchsepia},
    xlabel={$x$}, ylabel={$f(x)$},
    xmin=-2.6, xmax=2.6, ymin=-0.15, ymax=1.25,
    xtick={-2,-1,0,1,2}, ytick={0.5,1},
    tick label style={font=\small\sffamily, color=docink},
    label style={font=\small\sffamily, color=docink},
    samples=120, domain=-2.5:2.5,
  ]
    \addplot[draw=none, fill=benchlamp,
      domain=-2.5:2.5] {exp(-x^2/2)} \closedcycle;
    \addplot[thick, color=docaccent] {exp(-x^2/2)};
    \addplot[dashed, color=benchsepia,
      domain=-0.2:2.2]
      {exp(-0.5) - exp(-0.5)*(x-1)};
    \addplot[only marks, mark=*, mark size=2pt,
      color=benchink] coordinates {(1,{exp(-0.5)})};
    \node[anchor=south west, font=\small\sffamily,
      color=docink] at (axis cs:1.05,0.62)
      {tangente en $x=1$};
    \node[anchor=south, font=\small\sffamily,
      color=benchsepia] at (axis cs:1.35,0.95)
      {$f(x)=e^{-x^{2}/2}$};
  \end{axis}
\end{tikzpicture}
\end{code}

\subsection{Los preajustes}

El valor de \texttt{style} se fija una vez en \texttt{project.toml}, bajo
\texttt{[diagrams]}, y cada entorno puede cambiarlo con su propio atributo.
Los cuatro preajustes comparten el mismo criterio en Mermaid, Graphviz y D2:
\texttt{editorial} usa una paleta sobria con un solo acento para documentos a
color, \texttt{monochrome} dibuja solo en grises para la impresión en blanco y
negro, \texttt{technical} elimina rellenos y usa etiquetas de trazo uniforme
para esquemas, y \texttt{default} conserva el aspecto nativo de cada motor.
Omitir el atributo y la clave de proyecto equivale a \texttt{default}, de
modo que ningún documento existente cambia de aspecto.

  % Capitulo 4: matematicas con referencias cruzadas y derivacion.
\section{Matemáticas}\label{sec:matematicas}

La composición matemática es la de \LaTeX{} de siempre, con
\texttt{amsmath} cargado en el preámbulo. En línea, la fórmula
cuadrática $x = \frac{-b \pm \sqrt{b^2 - 4ac}}{2a}$ aparece dentro del
párrafo sin romper el ritmo, y la identidad de Euler $e^{i\pi} + 1 = 0$
recuerda que cinco constantes caben en una línea. La varianza de una
muestra, $\sigma^2 = \frac{1}{n}\sum_{i=1}^{n}(x_i - \mu)^2$, conecta este
capítulo con el guion de estadística del capítulo~\ref{sec:codigo}.

\subsection{Ecuaciones numeradas}

La integral de Gauss tiene valor exacto conocido, recogido en la
ecuación~\eqref{eq:gauss}. Su demostración clásica eleva la integral al
cuadrado y cambia a coordenadas polares; el resultado, $\sqrt{\pi}$,
reaparece cada vez que una normal necesita normalizarse.

\begin{equation}\label{eq:gauss}
  \int_{-\infty}^{\infty} e^{-x^2}\,dx = \sqrt{\pi}
\end{equation}

Junto a ella, la fórmula de Euler relaciona la exponencial compleja con el
seno y el coseno, como muestra la ecuación~\eqref{eq:euler}. Entre la
ecuación~\eqref{eq:gauss} y la ecuación~\eqref{eq:euler} queda
demostrado que las referencias cruzadas con \texttt{\textbackslash eqref}
funcionan: cada número entre paréntesis enlaza con su ecuación y
\texttt{texforge check} protesta si la etiqueta no existe.

\begin{equation}\label{eq:euler}
  e^{ix} = \cos x + i\,\sin x
\end{equation}

\subsection{Una derivación}

La media aritmética minimiza la suma de desvíos al cuadrado, y la
derivación cabe en un entorno \texttt{align} con una sola línea
numerada. Derivar respecto de $\mu$ e igualar a cero da directamente el
estimador, como se ve en la ecuación~\eqref{eq:derivacion}.

\begin{align}
  S(\mu) &= \sum_{i=1}^{n}(x_i - \mu)^2 \nonumber \\
  \frac{dS}{d\mu} &= -2\sum_{i=1}^{n}(x_i - \mu) \nonumber \\
  \hat{\mu} &= \frac{1}{n}\sum_{i=1}^{n}x_i \label{eq:derivacion}
\end{align}

La ecuación~\eqref{eq:derivacion} es la misma media que calcula la
función \texttt{mean} del listado~\ref{lst:archivo}: la teoría de este
capítulo y el código del capítulo~\ref{sec:codigo} cuentan la misma
historia por dos caminos, y ambos llegaron al PDF desde texto plano.

\subsection{Matrices}

Los vectores y las matrices se componen con los entornos de amsmath, y las
referencias funcionan igual que con las ecuaciones: la matriz identidad de
la ecuación~\eqref{eq:identidad} reaparece cada vez que un sistema lineal
necesita su elemento neutro.

\begin{equation}\label{eq:identidad}
  \mathbf{I} = \begin{pmatrix}
    1 & 0 & \cdots & 0 \\
    0 & 1 & \cdots & 0 \\
    \vdots & \vdots & \ddots & \vdots \\
    0 & 0 & \cdots & 1
  \end{pmatrix}
\end{equation}

La ecuación~\eqref{eq:identidad} multiplica sin cambiar, como el uno de la
aritmética: para toda matriz cuadrada vale que el producto con la identidad
devuelve la matriz intacta. Entre la integral de Gauss de la
ecuación~\eqref{eq:gauss} y esta identidad queda cubierto lo esencial:
integrales, exponenciales complejas, derivaciones alineadas y álgebra lineal
conviven en el mismo documento sin paquetes adicionales.

  % Capitulo 5: listados de codigo en los cuatro estilos, mas ink.
\section{Listados de código}\label{sec:codigo}

texforge resalta el código fuente durante la compilación: cada bloque
\texttt{code} se reescribe sobre la copia de compilación y solo necesita
\texttt{color.sty}, sin \texttt{listings} ni \texttt{minted}. El lenguaje se
indica con \texttt{lang} y la numeración con \texttt{numbers}. El tema de
color y el estilo por defecto se eligen una vez para todo el documento en
\texttt{project.toml}, aquí \texttt{one-light} con el estilo
\texttt{light}, y cada bloque puede pedir su propio estilo: hay
cuatro --- \texttt{light}, \texttt{light-mono}, \texttt{dark} y
\texttt{dark-mono} --- más el quinto estilo \texttt{ink}, en negro.

\subsection{Python}

El listado~\ref{lst:python} aproxima la integral de Gauss del
capítulo~\ref{sec:matematicas} por cuadratura: el mismo $\sqrt{\pi}$ de la
ecuación~\eqref{eq:gauss}, ahora calculado a fuerza bruta.

\begin{code}[lang=python, numbers=true, caption={Cuadratura de la integral de Gauss}, label={lst:python}]
import numpy as np

def integrar_gauss(n_puntos=1000):
    """Aproxima la integral de Gauss por cuadratura."""
    x = np.linspace(-10, 10, n_puntos)
    dx = x[1] - x[0]
    y = np.exp(-x**2)
    return np.sum(y) * dx

if __name__ == "__main__":
    print(f"Integral = {integrar_gauss():.6f}")
\end{code}

\subsection{Rust}

El listado~\ref{lst:fib} muestra la definición clásica de Fibonacci en
Rust con el estilo \texttt{light-mono}: sin color, las palabras clave van en
negrita y los comentarios en cursiva, que es lo que funciona impreso en
blanco y negro.

\begin{code}[lang=rust, style=light-mono, numbers=true, caption={Fibonacci en Rust}, label={lst:fib}]
fn fibonacci(n: u64) -> u64 {
    match n {
        0 => 0,
        1 => 1,
        _ => fibonacci(n - 1) + fibonacci(n - 2),
    }
}

fn main() {
    println!("fib(10) = {}", fibonacci(10));
}
\end{code}

\subsection{Shell}

El listado~\ref{lst:shell} usa el estilo \texttt{dark}, con el aire de una
terminal. Los caracteres especiales de \LaTeX{} se imprimen tal cual, porque
cada carácter se escapa antes de llegar al motor.

\begin{code}[lang=bash, style=dark, numbers=true, caption={Limpieza del directorio de compilación}, label={lst:shell}]
#!/usr/bin/env bash
set -euo pipefail

texforge build
texforge pdf info
texforge clean
\end{code}

\subsection{TOML}

El listado~\ref{lst:toml} es la configuración de resaltado de este
documento, con el estilo \texttt{dark-mono}: fondo oscuro y una sola tinta
clara, negritas y cursivas en lugar de colores.

\begin{code}[lang=toml, style=dark-mono, numbers=true, caption={Resaltado del propio documento}, label={lst:toml}]
[highlight]
theme = "one-light"
style = "light"
\end{code}

\subsection{Código desde un archivo}

El listado~\ref{lst:archivo} no está escrito en el \texttt{.tex}: se lee
del guion \texttt{scripts/stats.py} con \texttt{file} y un rango con
\texttt{lines}. Los números son los del archivo, y un archivo ausente o un
rango inválido fallan la compilación y \texttt{texforge check}.

\begin{code}[lang=python, file=scripts/stats.py, lines=5-24, numbers=true, caption={Funciones leídas del archivo}, label={lst:archivo}]
\end{code}

\subsection{Estilo ink}

El mismo archivo con \texttt{style=ink} usa solo negro: distingue los
elementos por peso y forma, con palabras clave y números en negrita,
nombres en negrita cursiva y comentarios en cursiva.

\begin{code}[lang=python, file=scripts/stats.py, lines=5-24, numbers=true, style=ink]
\end{code}

\begin{novabox}
  El estilo \texttt{ink} es nuevo en la 1.0: ignora el tema y compone en
  negro sobre blanco, pensado para imprimir o fotocopiar.
\end{novabox}

\subsection{Pies y etiquetas}

Todo bloque puede llevar pie y etiqueta: \texttt{caption=} escribe la línea
de Listado con su número y \texttt{label=} le da nombre para citarlo con
\texttt{\textbackslash ref}, como hacen los listados de este capítulo.
\texttt{pos=} decide si flota o se queda donde está escrito, y
\texttt{\textbackslash listoflistings} reúne todos los pies en un índice
propio, igual que el de figuras y tablas.

\subsection{Copiar sin números}

Los números de línea se omiten al copiar en los visores que respetan el
\texttt{ActualText} del PDF, y \texttt{numbers=false} los quita en todas
partes, como en el bloque siguiente.

\begin{code}[lang=python, file=scripts/stats.py, lines=5-24, numbers=false]
\end{code}

\begin{notebox}
  Pegar de un listado numerado no arrastra los números en los visores
  modernos: el PDF guarda el texto real aparte de lo que se ve. Si tu flujo
  exige páginas sin números, \texttt{numbers=false} los retira del todo.
\end{notebox}

  % Capitulo 6: bibliografia con seis citas distintas.
\section{Bibliografía}\label{sec:bib}

Las citas viven en \texttt{bib/references.bib} y se resuelven al componer,
con el estilo \texttt{plain} al final del documento. Este capítulo cita
seis entradas distintas para demostrar que el mecanismo funciona desde
cualquier capítulo: tres herramientas reales que sostienen el flujo y tres
entradas históricas del ejemplo original.

El motor que lo compone todo es tectonic~\cite{tectonic}, el corazón
embebido que presentó el capítulo~\ref{sec:que-es}. Los diagramas de
secuencia hablan Mermaid~\cite{mermaid} y los grafos dirigidos hablan el
lenguaje DOT de Graphviz~\cite{graphviz}; ambos se citaron ya en el
capítulo~\ref{sec:diagramas}, donde se componen.

Junto a ellas conviven las tres entradas originales del ejemplo: un
artículo de revista~\cite{example}, un libro~\cite{example2} y un recurso
en línea~\cite{example3}. Se conservan tal cual para que los documentos
que ya las citaban sigan compilando, y sirven aquí de testigos de que las
citas clásicas --- artículo, libro y web --- siguen resolviéndose como
siempre. \texttt{texforge check} avisa de toda clave \texttt{\textbackslash
  cite} sin entrada en el \texttt{.bib} antes de compilar.

El flujo bibliográfico no tiene misterio: cada obra vive una vez en
\texttt{bib/references.bib} con su clave, se cita en el texto con
\texttt{\textbackslash cite} y aparece formateada al final, ordenada según
el estilo elegido, aquí \texttt{plain}. Las claves conviene que sean
cortas y memorables --- el apellido, la herramienta, el año --- porque se
escriben a mano en cada cita. Y cuando el documento crece, el mismo
\texttt{.bib} sirve a varios capítulos sin duplicar nada: este capítulo
cita seis entradas distintas repartidas por todo el documento, y la lista
final las reúne a todas sin que nadie las ordene a mano.

\begin{notebox}
  Las tres primeras entradas vienen del ejemplo original y se conservan a
  propósito: los documentos que ya las citaban siguen compilando sin
  cambios, y aquí demuestran que lo clásico convive con lo nuevo.
\end{notebox}

  % Capitulo 7: novedades de la version 1.0, tabla y lista.
\section{Novedades de la 1.0}\label{sec:novedades}

La versión 1.0 redondea el flujo con comandos que antes no existían y con
detalles de acabado en los listados. Hasta ahora crear un proyecto con
datos era manual, las plantillas propias no existían, el entorno se
diagnosticaba a ojo, actualizar era descargar a mano, la estructura solo
se veía en el PDF y los listados vestían un único estilo con números que
se pegaban al copiar. La tabla~\ref{tab:novedades} recoge lo que cambia,
con su documentación de referencia; la lista que sigue cuenta lo que
cambia en el día a día.

\begin{table}[H]
  \centering
  \caption{Novedades de la versión 1.0 y dónde leer más}
  \label{tab:novedades}
  \begin{tabularx}{\linewidth}{@{}l>{\raggedright\arraybackslash}X>{\raggedright\arraybackslash}X@{}}
    \toprule
    \sffamily\small\bfseries Novedad & \sffamily\small\bfseries Qué hace & \sffamily\small\bfseries Dónde leer más \\
    \midrule
    \texttt{new -{}-values} & Rellena la plantilla desde un fichero de valores y lo guarda como \texttt{values.toml}. & Referencia de \texttt{new} \\
    \texttt{template add} & Registra una carpeta local como plantilla. & Referencia de \texttt{template} \\
    \texttt{skill} & Expone la destreza de agente del binario. & Referencia de \texttt{skill} \\
    \texttt{doctor} & Diagnostica y actualiza el Tectonic gestionado. & Referencia de \texttt{doctor} \\
    \texttt{update} & Trae el binario nuevo con la firma verificada. & Referencia de \texttt{update} \\
    \texttt{outline} & Dibuja el árbol de secciones. & Referencia de \texttt{outline} \\
    \texttt{pdf} & Inspecciona el PDF: texto, fuentes, páginas y chequeo. & Referencia de \texttt{pdf} \\
    Estilo \texttt{ink} & Listados en negro, sin depender del tema. & Listados de código \\
    DejaVu por defecto & Los listados usan DejaVu Sans Mono salvo decisión propia. & Listados de código \\
    Números que no se copian & El \texttt{ActualText} deja los números fuera del portapapeles. & Listados de código \\
    \bottomrule
  \end{tabularx}
\end{table}

\begin{novabox}
  Todo este documento es nuevo en la 1.0 por dentro: se recompuso como pieza
  editorial sobre el lenguaje visual de UniverLab, con portada, cabecera,
  cajas y paleta propias, en lugar del artículo neutro anterior.
\end{novabox}

\begin{itemize}
  \item Crear con valores: \texttt{texforge new tesis -t tesis -{}-values
    datos.toml} deja cada proyecto con sus datos desde el primer minuto, y
  el fichero usado queda archivado como \texttt{values.toml}.
  \item Plantillas propias: cualquier carpeta con un \texttt{main.tex} y su
  \texttt{project.toml} se vuelve plantilla con \texttt{template add}, sin
  pasar por el registro remoto.
  \item Entorno sano: \texttt{doctor} dice qué versión del motor hay,
  cuánto ocupa la caché y qué diccionarios están instalados, y trae
  el motor gestionado al día sin que tengas que pedirlo.
  \item Actualización confiada: \texttt{update} comprueba la firma antes
  de sustituir el binario, y \texttt{update -{}-check} solo pregunta si hay
  algo nuevo.
  \item Leer el PDF sin abrirlo: \texttt{pdf text} para la accesibilidad,
  \texttt{pdf info} para las fuentes, \texttt{pdf pages} para comparar
  compilaciones y \texttt{pdf check} para cazar palabras perdidas.
  \item Listados serios: pie con numeración, etiqueta para citarlo,
  lectura desde archivo con \texttt{file} y \texttt{lines}, y el estilo
  \texttt{ink} cuando el destino es el papel.
\end{itemize}

  % Capitulo 8: listas con proposito.
\section{Listas}\label{sec:listas}

Tres entornos bastan para casi todo lo que hay que enumerar, y este capítulo los
usa a los tres en el orden en que conviene aprenderlos: primero lo que no
tiene orden, luego lo que sí lo tiene y por último lo que necesita nombre.
Cada lista tiene su oficio. La viñeta enumera, la numerada ordena y la de
descripción define; usar cada una para lo suyo es lo que separa un
documento legible de un amasijo de puntos.

Lo que texforge evita instalar cabe en una viñeta, porque el orden no
importa y la lista puede crecer sin reescribir nada:

\begin{itemize}
  \item ni TeX Live ni MiKTeX;
  \item ni Inkscape ni Node para los diagramas;
  \item ni \texttt{minted} ni \texttt{listings} para el código.
\end{itemize}

El ciclo de trabajo, en cambio, pide números, porque el orden sí
importa y cada paso presupone el anterior:

\begin{enumerate}
  \item escribir el fuente en texto plano;
  \item revisar con \texttt{texforge check};
  \item formatear con \texttt{texforge fmt};
  \item componer con \texttt{texforge build};
  \item inspeccionar con \texttt{texforge pdf}.
\end{enumerate}

Y los conceptos con nombre propio piden una lista de descripción, donde el
término manda y la definición sigue:

\begin{description}
  \item[Preámbulo] la cabecera del \texttt{.tex}: clase, paquetes,
  paleta, cabeceras y cajas. Se escribe una vez y gobierna todo el
  documento.
  \item[Entorno] un bloque con sintaxis propia, como \texttt{mermaid} o
  \texttt{code}: texforge lo detecta al compilar y lo dibuja.
  \item[Etiqueta] el nombre que \texttt{\textbackslash label} o
  \texttt{label=} da a una sección, figura, tabla, ecuación o
  listado para citarlo con \texttt{\textbackslash ref}.
\end{description}

\begin{notebox}
  La regla práctica: si reordenar los puntos cambia el sentido, numera;
  si cada punto necesita su término, describe; en lo demás, viñetas.
\end{notebox}

  % Capitulo 9: conclusion.
\section{Conclusión}\label{sec:conclusion}

Este documento empezó como un artículo neutro que enseñaba piezas
sueltas y termina como una pieza editorial que las usa todas a la vez. La
portada con su banda granate, el resumen de un vistazo, el índice, las
nueve secciones, las dos tablas, los cuatro diagramas, las ecuaciones con
sus referencias, los seis listados y las seis citas salieron de una sola
orden: \texttt{texforge build}. Nada se pegó como imagen y nada se
compuso fuera del flujo que el propio documento describe.

texforge demuestra aquí lo que promete en su manual: escribir en texto
plano, revisar sin compilar, formatear sin pensar y componer sin una
distribución externa. Los diagramas viajan con el fuente, el código se
resalta solo, la ortografía se revisa en el idioma del documento y el PDF
resultante se inspecciona desde la terminal. Quien publique con estas
piezas --- una tesis, un informe, un manual --- hereda el mismo andamiaje:
una paleta, dos familias tipográficas y un preámbulo que ya resolvió
las decisiones visuales.

Queda trabajo por delante, como en todo proyecto vivo, pero el punto de
partida es este: un documento real, de presentación cuidada, producido
íntegramente con la herramienta que describe. Quien lo compila obtiene lo
mismo en cualquier máquina, porque el motor viaja con la herramienta y el
fuente no depende de nada instalado a mano: el binario y el texto bastan.
Si algo de lo que ves te resulta útil, su fuente está en
\texttt{examples/texforge-capabilities} esperando tu próxima compilación,
y el manual sigue en \texttt{docs/cli-reference.md}.


  % La última entrada cabe en la página anterior con una línea de más.
  \enlargethispage{\baselineskip}
  \bibliographystyle{plain}
  \bibliography{bib/references}

\end{document}
